
Lúc 10h sáng 13/5/2018 GS. TSKH Phan Đình Diệu đã thanh thản về an nghỉ nơi Vĩnh hằng hưởng thọ 82 tuổi. Xin gửi lời chia buồn sâu sắc  đến toàn thể gia quyến của Giáo sư.

Đã nhiều báo có bài viết về nhà Khoa học lớn của đất nước trong lĩnh vực Toán học, Triết học...Phan Đình Diệu.Tôi nhớ về GS Phan Đình Diệu như Người anh cả của ngành Công nghệ thông tin (CNTT) và khoa học tính toán của nước nhà. Ngày hôm nay chúng ta được hưởng những thành quả ứng dụng CNTT và Internet ở Việt Nam là có sự góp sức từ thuở ban đầu không nhỏ của anh khi triển khai Nghị Quyết 49/CP “ về phát triển Công nghệ thông tin của nước ta trong những năm 90” do Thủ tướng Võ Văn Kiệt ký ban hành ngày 4/8/1993. Anh cũng là Phó trưởng ban chuyên trách, Thường trực Ban chỉ đạo Chương trình Quốc gia đầu tiên của Việt Nam ( 1994-1998). 

Anh chính là người sáng lập và là Chủ tịch đầu tiên của Hội Tin học Việt Nam ( 1988) là Viện trưởng đầu tiên cuả Viện Khoa học Tính toán và Điều khiển, nay là Viện Công nghệ Thông tin ( 1977).

Vĩnh biệt người Thủ trưởng, người Thầy và  người Anh đã  làm việc cùng tôi, dìu dắt tôi từ tuổi thanh niên đến suốt 20 năm (1977-1997). Kính chúc hương hồn anh an giấc ngàn thu trên cõi Niết Bàn. Mong anh luôn phù hộ cho sự nghiệp phát triển CNTT của nước nhà vượt qua mọi thăng trầm để sánh ngang với bạn bè Năm Châu. 

% Kính tặng Anh 4 câu Thơ như một lời tri ân:

% TIẾC THƯƠNG ANH

% Bầu trời lại tắt một ngôi sao

% Bạn hữu xa anh nước mắt trào

% Trí tuệ uyên thâm đời bác học

% Đức tài trọn vẹn sống thanh cao

% -------------

%Chú gửi cháu 2 bài Thơ Đường luật trọn dạng “ Thất ngôn bát cú”
%kính tặng hương hồn Bố cháu. Sinh thời bố cháu là một người rất uyên bác về Thơ Đường đấy. 
Hai bài Thơ kính tặng hương hồn  người Thầy, người Anh Phan Đình Diệu:

TIẾC THƯƠNG ANH

Bầu trời lại khuyết một ngôi sao

Bạn hữu xa anh nước mắt trào

Trí tuệ uyên thâm đời bác học

Đức tài trọn vẹn sống thanh cao

Lo về quốc pháp lời trang trọng

Nghĩ tới gia phong nét tự hào

Khí phách sĩ phu Nhân Nghĩa Dũng

Vui cùng thơ phú thú tiêu dao

(15/5/2018)

VĨNH BIỆT ANH 

Tiễn biệt anh rồi dạ ngẩn ngơ

Bồng lai Tiên cảnh đón vô bờ

Nơi kia Anh nghỉ ngàn thu mộng

Cõi ấy Anh nằm vạn giấc mơ

Quốc sách thông tin còn để lại

Lời khuyên Giáo dục vẫn đang chờ

Thiên Đường dạo ngắm hồn thanh thản

Vẹn nghĩa chung tình đẹp ý Thơ

(18/5/2018)

/Nguyễn Công Hoá/